\documentclass[10pt,a4paper]{amsart}
\usepackage[margin=19mm]{geometry}
\usepackage{amsmath,amssymb,mathtools}
\usepackage[scr=rsfs,cal=euler]{mathalfa}
\usepackage{xfrac}
\usepackage[unicode,hidelinks,bookmarks=false]{hyperref}
\newcommand{\ie}{i.e.\ }
\newcommand{\cf}{cf.\ }
\newcommand{\resp}{respectively\ }
\newcommand{\Subsection}{\subsection}
\providecommand{\nobreakdash}{\nobreak\hbox{-}}
\setlength{\emergencystretch}{2em}
\multlinegap0pt
\usepackage{bm}
\usepackage{bbm}
\usepackage{dsfont}
\usepackage{xspace}
\usepackage{cancel}
\usepackage{mathtools}
\usepackage{amsthm}
\makeatletter
\@ifundefined{thm}{\newtheorem{thm}{Theorem}}{}
\makeatother
\title[Wilf Equivalence for Length-Three Patterns and Flat POPs, an]{Wilf Equivalence for Length-Three Patterns and Flat POPs, and a Conjecture of Qiu and Remmel}
\author{Shiqi Cao, Sergey Kitaev,, Yuxin Wu}
\date{Source: arXiv:2608.18461v1 (CC BY 4.0)}
\begin{document}
\maketitle

\noindent\emph{Source and licence.} Wilf Equivalence for Length-Three Patterns and Flat POPs, and a Conjecture of Qiu and Remmel. Version arXiv:2608.18461v1 (CC BY 4.0), distributed under the Creative Commons Attribution 4.0 International licence, as recorded in the arXiv licence metadata for this exact version and shown on the abstract page https://arxiv.org/abs/2608.18461v1. Full rights notice: licenses/arxiv-2608.18461-CC-BY-4.0.txt.\par\smallskip

\noindent\textit{Editorial scope.} Counted unit: the Thm whose source label is \texttt{thm:Phi-Psi-preserve-POP} in the article (source0.tex lines 939--952, source label \texttt{thm:Phi-Psi-preserve-POP}), delivered together with the article's own complete original proof of it (source0.tex lines 954--1095, one \texttt{proof} environment of 142 lines). No mathematics is written, repaired or re-derived: statement and proof are the article's own text line for line. Required context retained uncounted: none. Every definition, display and label of those lines is retained unchanged. Omitted from this package and not redistributed: 1--938, 953--953, 1096--2290. Nothing inside a retained excerpt is omitted or altered: every retained body is delivered exactly as the article's lines are written, so no omission receipt of any kind accompanies this package. Citations: the retained excerpts carry no citation command at all, so no bibliography entry of the article is transcribed and no reference is redistributed. Reference adaptation: none was needed and none was made; every reference inside the delivered unit resolves inside it, so no unresolved reference or citation is produced. Printed slips kept literal: no printed word, symbol, hypothesis or number of the counted statement or of its proof was altered. Layout and class: the article's own class is \texttt{article} with packages including amsmath, amsthm, amsfonts, amssymb, latexsym, diagbox, mathtools, array, subcaption, booktabs, tkz-graph, graphics, graphicx, multirow; the wrapper is the lane's pinned \texttt{amsart} editorial wrapper at 10pt on A4, which restates the theorem-like environments the retained text needs and adds the article's own macro definitions that the retained text needs (none), with extra wrapper packages (bm, bbm, dsfont, xspace, cancel, mathtools). The delivered numbering is the unit's own. Assets: no retained span contains a figure environment, a graphics-inclusion command, a PDF-page inclusion, a table or a TikZ picture; no image file is copied into this package, the package declares no asset dependency and no image file is redistributed with it. Licence: version arXiv:2608.18461v1 of this article is distributed under the Creative Commons Attribution 4.0 International licence, as recorded in the arXiv licence metadata for this exact version and shown on the abstract page https://arxiv.org/abs/2608.18461v1. A full rights notice is delivered at licenses/arxiv-2608.18461-CC-BY-4.0.txt. Characters: every retained excerpt is pure ASCII, so the delivered engine, XeLaTeX with its default Latin Modern text font, reports no missing character.
\begin{thm}\label{thm:Phi-Psi-preserve-POP}
For $2\leq x\leq\ell$, the maps $\Phi$ and $\Psi$ restrict to
mutually inverse bijections
$\Phi:
\operatorname{Av}_n(321,P_{\ell,x})
\longrightarrow
\operatorname{Av}_n(231,P_{\ell,x})$
and
$
\Psi:
\operatorname{Av}_n(231,P_{\ell,x})
\longrightarrow
\operatorname{Av}_n(321,P_{\ell,x})$.
\end{thm}

\begin{proof}
Let $\pi\in\operatorname{Av}_n(321,P_{\ell,x})$ and
$\sigma\in\operatorname{Av}_n(231,P_{\ell,x})$.
For convenience, define
$
\pi^{\langle a\rangle}
=
\bigl(
\varphi^{\langle a\rangle}\circ\cdots\circ
\varphi^{\langle1\rangle}
\bigr)(\pi)$
and
$
\sigma^{\langle a\rangle}
=
\bigl(
\psi^{\langle a\rangle}\circ\cdots\circ
\psi^{\langle1\rangle}
\bigr)(\sigma)$.
Recall that an entry $v$ can occupy the $x$-th position of an occurrence
of $P_{\ell,x}$ if and only if
$Q_{\mathrm{II}}(v)\geq x-1$ and
$Q_{\mathrm I}(v)\geq\ell-x$.
We first prove that $\Phi$ preserves avoidance of $P_{\ell,x}$. We
show inductively that if an occurrence of $P_{\ell,x}$ appears in
$\pi^{\langle a\rangle}$, then none of $1,2,\ldots,a+1$
can occupy its $x$-th position.

First consider the entry $1$. It is easy to check that rearranging entries does not change
the numbers of entries in the first and second quadrants of $1$.
Therefore,
$Q_{\mathrm I}^{\pi^{\langle a\rangle}}(1)
=
Q_{\mathrm I}^{\pi}(1)$ and
$Q_{\mathrm{II}}^{\pi^{\langle a\rangle}}(1)
=
Q_{\mathrm{II}}^{\pi}(1)$.
Since $\pi$ avoids $P_{\ell,x}$, the entry $1$ cannot occupy the
$x$-th position of an occurrence of $P_{\ell,x}$ in any
$\pi^{\langle a\rangle}$.

Now suppose inductively that, in $\pi^{\langle a-1\rangle}$, none of
$1,\ldots,a$ can occupy the $x$-th position of an occurrence. The
permutations $\pi^{\langle a-1\rangle}$ and
$\pi^{\langle a\rangle}$ differ only by
$\varphi^{\langle a\rangle}$, which rearranges only entries larger than
$a$. Hence, for every $j\leq a$,
$Q_{\mathrm I}^{\pi^{\langle a\rangle}}(j)
=
Q_{\mathrm I}^{\pi^{\langle a-1\rangle}}(j)$
and
$Q_{\mathrm{II}}^{\pi^{\langle a\rangle}}(j)
=
Q_{\mathrm{II}}^{\pi^{\langle a-1\rangle}}(j)$.
Thus none of $1,\ldots,a$ can occupy the $x$-th position in
$\pi^{\langle a\rangle}$. It remains to consider $a+1$.
We distinguish two cases.

\noindent \textbf{Case 1:} $a+1$ has not been acted on by any of
$\varphi^{\langle1\rangle},\ldots,\varphi^{\langle a\rangle}$.

In this case, whenever $\varphi^{\langle i\rangle}$ is applied, the entry
$i$ lies to the left of $a+1$. 
Thus all the rearrangements take place
to the left of $a+1$, and hence
$Q_{\mathrm I}^{\pi^{\langle a\rangle}}(a+1)
=
Q_{\mathrm I}^{\pi}(a+1)$
and
$Q_{\mathrm{II}}^{\pi^{\langle a\rangle}}(a+1)
=
Q_{\mathrm{II}}^{\pi}(a+1)$.
So $a+1$ cannot occupy the $x$-th position in $\pi^{\langle a\rangle}$.

\noindent \textbf{Case 2:} $a+1$ has been acted on by some
$\varphi^{\langle i\rangle}$.

Let $\varphi^{\langle i\rangle}$ be the last operation acting on $a+1$.
Suppose, for contradiction, that $a+1$ occupies the $x$-th position
of an occurrence of $P_{\ell,x}$. Then
$Q_{\mathrm{II}}^{\pi^{\langle a\rangle}}(a+1)\geq x-1$ and
$Q_{\mathrm I}^{\pi^{\langle a\rangle}}(a+1)\geq\ell-x$.
So we have 
$
Q_{\mathrm{II}}^{\pi^{\langle a\rangle}}(i)
\geq
Q_{\mathrm{II}}^{\pi^{\langle a\rangle}}(a+1)+1
\geq x$.
Since $i\leq a$, the induction hypothesis implies that $i$ cannot
occupy the $x$-th position of an occurrence. Therefore,
$Q_{\mathrm I}^{\pi^{\langle a\rangle}}(i)<\ell-x$.
If $x=\ell$, the last inequality is already impossible. Otherwise,
since
$Q_{\mathrm I}^{\pi^{\langle a\rangle}}(a+1)\geq\ell-x$ and
$Q_{\mathrm I}^{\pi^{\langle a\rangle}}(i)<\ell-x$,
there exists an entry $j>a+1$ such that
$\operatorname{Pos}_{\pi^{\langle a\rangle}}(a+1)
<
\operatorname{Pos}_{\pi^{\langle a\rangle}}(j)
<
\operatorname{Pos}_{\pi^{\langle a\rangle}}(i)$.
This contradicts the fact that the
second quadrant subword of $i$ is decreasing. 
Hence $a+1$ cannot
occupy the $x$-th position of an occurrence of $P_{\ell,x}$.
This completes the induction. Therefore, $\Phi(\pi)\in\operatorname{Av}_n(231,P_{\ell,x})$.

We next prove that $\Psi$ preserves avoidance of $P_{\ell,x}$. 
We proceed with the same line of reasoning to prove it.

The position of $1$ does not change, and neither do its
first and second quadrant counts. 
Now suppose the assertion holds
for $\sigma^{\langle a-1\rangle}$. Since
$\psi^{\langle a\rangle}$ rearranges only entries larger than $a$,
it remains only to consider $a+1$.
Again, there are two cases.

\noindent \textbf{Case 1:} $a+1$ has not been acted on by any
$\psi^{\langle i\rangle}$.

In this case, the numbers of entries in the first and second quadrants
of $a+1$ are unchanged. Since $\sigma$ avoids $P_{\ell,x}$, the entry
$a+1$ cannot occupy the $x$-th position of an occurrence in
$\sigma^{\langle a\rangle}$.

\noindent \textbf{Case 2:} $a+1$ has been acted on by some
$\psi^{\langle i\rangle}$.

Let $\psi^{\langle i\rangle}$ be the last operation acting on $a+1$.
This operation arranges the second-quadrant subword of $i$ in
increasing order. Consequently, there is no entry larger than $a+1$
to the left of $a+1$. Hence
$Q_{\mathrm{II}}^{\sigma^{\langle a\rangle}}(a+1)=0$.
Since $x\geq2$, we have
$Q_{\mathrm{II}}^{\sigma^{\langle a\rangle}}(a+1)
=0<x-1$.
Therefore $a+1$ cannot occupy the $x$-th position of an occurrence of
$P_{\ell,x}$.
The induction is complete, and hence
$\Psi(\sigma)\in\operatorname{Av}_n(321,P_{\ell,x})$.
\end{proof}

\end{document}
